%!TEX root = /Users/markelikalderon/Documents/Git/perceptual_self-consciousness/perceptual_self-consciousness.tex
\chapter{Second Excurus} % (fold)
\label{cha:second_excurus}

\section{First-Person Perspective and the Phenomenal} % (fold)
\label{sec:first-person_perspective_and_the_phenomenal}

We have seen, in \emph{De somno}, perceptual self-consciousness and apperceptive unity are linked in that they admit of a common explanation. It is the common power (\emph{koinē dunamis}) that explains that we perceive that we see and hear and that we can discriminate white from sweet. A question arises why these seemingly distinct phenomena, perceptual self-consciousness and apperceptive unity, might be linked? We can ask this question both of Aristotle's account in \emph{De somno} but also of ourselves.

Asked of the \emph{De somno} account, the answer seems reasonably clear, at least if we accept that account on its own terms. When the animal sleeps, its power to perceive that it sees and its power to discriminate white from sweet are incapacitated. When one is asleep one is not aware of seeing the surrounding darkness, and not just because one is not seeing, one is not aware of not seeing. What would furnish one with this latter awareness is sight itself which is temporarily incapcitated in slumber. Nor can one discriminate sensibles from distinct \emph{genera}, and not just because one is not perceiving through the the relevant senses, but because the power to discriminate them is incapacitated in slumber. It is because the same power is incapacitated, the \emph{koinē dunamis}, when its corporeal instrument is at rest, the primary sense organ (presumably the heart, given Aristotle's cardiocentrism) and because this power is multi-track, being manifest both in perceptual self-consciousness and apperceptive unity, that these distinct phenomena are related.

If we ask the same question of ourselves, rather than of Aristotle's \emph{De somno} account, a deeper answer emerges, and one that may be hovering in the background of Aristotle's own reflection in \emph{De somno} and other works. So why might perceptual self-consciousness and apperceptive unity be linked in the way Aristotle suggests? On the interpretation advanced in the last chapter, perceptual self-consciousness is reflexive in character rather than higher-order. Here, then, is the suggestion: Reflexivity and unity are each essential components of first-person perspective.

\textsc{The Regress Argument}, if it is not simply invalid \citep[202]{Corkum:2010fa}, must involve the assumption that every perception is conscious as Themistius (\emph{In libros Aristotelis De anima paraphrasis}
83,15–6) suggests and that the avowed Cartesian, Sartre (\emph{L'Être et le Néant}), insists upon. If consciousness at least involves being aware from within that one perceives, and such awareness is reflexive, then having a conscious first-person perspective essentially invovles a reflexive component. Even if, like Leibniz (\emph{Principes de la Nature et de la Grâce} §4, \emph{La Monadologie} §14), one insists on distinguishing perception and apperception thus countenancing unconscious perception, one might still think that it is essential to the first-person perspective that one might apperceive. The difference is that the Cartesian position insists that one necessarily enjoys reflexive awareness whenever one perceives, whereas the Leibnizian position insists that necessarily one is capable of reflexive awareness, whether or not this power is exercised.

A conscious first-person perspective is a unified perspective. It is because one's first person perspective is unified that one can discriminate white from sweet. Without the requisite unity, the senses would be like so many separate Achaeans hiding in the Wooden Horse (\emph{Theaetetus} 184c–d; chapter~\ref{sub:the_argument_from_residual_imagery}). One sense may perceive white and another sweet, but neither individually nor in concert will be able to discriminate white from sweet. Such discriminatory activity is only possible from a unified perspective. It is not just our ability to discriminate sensibles from distinct \emph{genera} that would be compromised, but much more besides. Consider Augustine's description of responding to a bodily injury:
\begin{quote}
	When, for instance, there is an ache in the foot, the eye looks at it, the mouth speaks of it, and the hand reaches for it. (Augustine, \emph{De immortalitate animae} 16.25; \citealt[46]{Schopp:1947df})
\end{quote}
The coordinated activities of the eye, mouth, and hand simultaneously directed at the injured foot would only be possible from the unified perspective afforded by the soul. 

The first-person perspective is relevant as well to McCabe's \citeyearpar{McCabe:2007jb} criticism of \citet{Kosman:2014ab}. In a section ominously entitled \emph{The Snares of Consciousness}, McCabe puruses the charge of anachronism against Kosman. Kosman, had, by McCabe's lights, misattributed to Aristotle an account of phenomenal consciousness but no such conception is alleged to be found in the Aristotelian corpus. This, according to McCabe, is a modern intrusion that obscures more than it enlightens.

The terms of the charge are surprising, however, since Kosman never speaks of phenomenal consciousness. 

%section first-person_perspective_and_the_phenomenal (end)

\section{The Revenge Regress} % (fold)
\label{sec:the_revenge_regress}


% section the_revenge_regress (end)

\section{\emph{Sentio Ergo Sum}} % (fold)
\label{sec:sensato_ergo_sum}


% section sensato_ergo_sum (end)

% chapter second_excurus (end)